% Appendix A20 — cost-benefit deployment analysis (methodology).
% Source: ACCEPTANCE_CRITERIA.md L13 (cost-benefit: missed-melanoma / retake / referral cost,
% DALY/QALY framing) + STORY_FRAMEWORK A20. METHODOLOGY ONLY; dollar/DALY figures and the
% simulated operating curves are PENDING (M3) and not reported here (red line 4).
% Anonymized; clinical framing per R3 (simulation/decision-curve language, no reader claims).

\section{Cost--Benefit Deployment Analysis}\label{app:a20}

\subsection{Motivation}\label{app:a20-motiv}
A quality-triage system trades one cost (retakes and specialist referrals) against another (missed or
mis-staged lesions). Whether this trade is favourable cannot be read off accuracy alone; it depends on
the relative costs of each error and action. This appendix specifies the decision-analytic model used
to assess that trade; the numerical instantiation is reported in \cref{sec:exp-dca}.

\subsection{Cost Model}\label{app:a20-model}
For a case routed to action $a\in\{\actd,\acte,\actq,\actr\}$ with true label $y$ and prediction
$\hat{y}$, we assign an expected cost
\begin{equation}
C(a,y,\hat{y}) = c_{\mathrm{err}}(y,\hat{y}) + c_{\mathrm{act}}(a),
\label{eq:a20-cost}
\end{equation}
where $c_{\mathrm{err}}$ is the diagnostic-error cost (a false negative on a malignant lesion
dominates, reflecting downstream morbidity, optionally expressed in DALY/QALY units) and
$c_{\mathrm{act}}$ is the operational cost of the action (negligible for $\actd$; a capture cost for
$\actq$; a specialist-time cost for $\actr$; an inference/enhancement cost for $\acte$). All cost
constants are exposed as parameters and swept rather than fixed to a single jurisdiction.

\subsection{Net Benefit and Decision Curves}\label{app:a20-netbenefit}
We report \emph{net benefit} as a function of the decision threshold $p_t$,
\begin{equation}
\mathrm{NB}(p_t) = \frac{\mathrm{TP}}{N} - \frac{\mathrm{FP}}{N}\cdot\frac{p_t}{1-p_t},
\label{eq:a20-nb}
\end{equation}
and compare the closed-loop agent against (i) treat-all, (ii) treat-none, and (iii) direct diagnosis
without triage. Decision-curve dominance over a clinically relevant $p_t$ range is the headline
deployment claim; it requires no human readers and is computed entirely from held-out predictions.

\subsection{Triage Simulation}\label{app:a20-sim}
On top of the static cost model we simulate the closed loop: degraded inputs are routed by the policy
of Theorem~\ref{thm:agent}, retakes draw a higher-quality re-capture from the paired data,
enhancement applies $\Tw$, and referrals remove the case from automated decision. Sweeping the
quality thresholds $(\tauenh,\tauhigh)$ traces an operating curve of expected cost versus referral
rate, identifying the band where triage is economically favourable. The simulation inherits the
safety boundary of \cref{sec:exp-salvage} (low salvage rate on severely degraded inputs is intended,
not a failure).

\subsection{Reporting and Limitations}\label{app:a20-limits}
We report the cost model as a \emph{sensitivity analysis} over the cost constants rather than a single
point estimate, since real costs vary by health system. The analysis is a deployment-feasibility
argument, not a health-economic evaluation; a formal cost-effectiveness study with jurisdiction-specific
inputs is out of scope and listed as future work in \cref{sec:discussion}.
